%% ===================================================================
%% Signal-based decomposition of battery check-up tests
%% Target: Journal of Power Sources / Journal of Energy Storage
%% Class:  elsarticle (Elsevier). Install via TeX Live `elsarticle`,
%%         or start from the Overleaf "Elsevier article" template.
%% Build:  latexmk -pdf main.tex
%% ===================================================================
\documentclass[preprint,3p,number]{elsarticle}

\usepackage{amsmath}
\usepackage{booktabs}
\usepackage{graphicx}
\usepackage{siunitx}
\usepackage{xcolor}
\usepackage{url}

\sisetup{detect-all, group-separator={\,}}

%% ---- draft scaffolding -------------------------------------------
%% Every unrun experiment is wrapped in \todobox so it is impossible to
%% mistake a placeholder for a result. Set \drafttrue to hide them all
%% before circulating externally.
\newif\ifdraft\drafttrue
\newcommand{\todobox}[2]{%
  \ifdraft
  \begin{center}
  \fcolorbox{red!60!black}{red!4}{%
    \begin{minipage}{0.94\linewidth}
    \small\textbf{\textcolor{red!60!black}{TODO --- #1}}\\[2pt]#2
    \end{minipage}}
  \end{center}
  \fi}
\newcommand{\todo}[1]{\ifdraft\textcolor{red!60!black}{[\textbf{TODO:} #1]}\fi}

\journal{Journal of Power Sources}

\begin{document}

\begin{frontmatter}

\title{Signal-based decomposition of battery check-up tests without a
step-by-step protocol description: chemistry-portable identification of
capacity, pulse and quasi-OCV segments across NMC, LFP and sodium-ion cells}

\author[isea]{\todo{author list}}
\address[isea]{\todo{affiliation}}

\begin{abstract}
Reference performance tests (RPTs, or check-ups) are the measurements that make
an ageing dataset interpretable: they contain the capacity test, the resistance
pulses and the quasi-open-circuit-voltage sweep from which state of health,
internal resistance and the OCV curve are derived. Yet the software that
processes public ageing datasets routinely discards them, because locating them
requires knowing the test protocol. Existing tools resolve this by trusting
metadata: either the cycler's own cycle index, or an experimenter-authored
step-index map that assigns a meaning to every step of every protocol. Both tie
the analysis to one laboratory's bookkeeping.
%
We present a method that locates and decomposes check-up tests from the
measured signal, requiring only a small set of physical cell and test
parameters --- voltage limits, nominal capacity and the target C-rates --- and
neither a step-by-step protocol description nor per-source code. The method
segments a raw time series into procedures, describes each segment with
scale-normalised features, and assigns capacity (CAP), pulse (PUL) and
quasi-OCV (qOCV) labels.
%
We report two experiments. First, an ablation over the segmentation stage shows
that the protocol columns it currently reads are not load-bearing: the cycler's
state column is redundant with the sign of the measured current to
\SI{0.0008}{\percent} of \num{131442261} samples, and boundary placement driven
by current alone recovers \SI{100}{\percent} of check-up segments intact --- no
segment merged or lost --- on nickel-manganese-cobalt (NMC), lithium iron
phosphate (LFP) and sodium-ion (SIB) cells. Second, a leave-one-chemistry-out
transfer test shows that a classifier trained on NMC alone recovers every
capacity test on unseen LFP and sodium-ion cells given only their physical
parameters, provided its label space is scale-free; a classifier trained on
absolute C-rates fails completely on LFP.
%
\todo{one closing sentence once items 3--5 have run}
\end{abstract}

\begin{keyword}
lithium-ion batteries \sep sodium-ion batteries \sep reference performance test
\sep check-up \sep unsupervised segmentation \sep data pipelines
\end{keyword}

\end{frontmatter}

%% ===================================================================
\section{Introduction}
\label{sec:intro}

Battery ageing studies produce two kinds of data. The bulk of it is cycling:
long, repetitive load profiles that age the cell. Interleaved with it, at
intervals of days to months, sits the \emph{check-up} or reference performance
test (RPT) --- a short, carefully specified sequence that measures the cell
rather than ageing it. A typical check-up contains a full capacity test at a
defined C-rate, a series of current pulses for resistance, and a low-rate
quasi-OCV sweep. Almost every quantity a battery paper reports about degradation
--- state of health, resistance growth, the shift of the OCV curve, differential
voltage features --- is extracted from those few hours, not from the cycling
around them.

Locating the check-up in a raw cycler export is therefore the precondition for
almost any downstream analysis, and it is exactly the step that current tooling
externalises to metadata. The recently published BatteryLife benchmark
\cite{batterylife} unifies sixteen ageing datasets, and does so with one
preprocessing script per source; the check-ups are not analysed but
\emph{removed}, in one case by a hard-coded current threshold, in another by
parsing only the files whose names mark them as cycling. PyProBE
\cite{pyprobe}, published in a software venue, takes the opposite and more
principled route --- it models the experiment explicitly --- but requires the
experimenter to author a step-index map (``steps 9--12 = discharge pulses'')
for every protocol, including manual reconciliation of cycler quirks.

These are not competing implementations of the same idea; they occupy different
points on a single axis, namely \emph{how much protocol metadata is needed to
find structure in the signal}. BatteryLife trusts the cycler's cycle index.
PyProBE trusts a human-authored step index. Both are reasonable engineering
choices, and both mean the analysis does not transfer to a source that indexes
its steps differently, or at all.

This paper asks whether the structure can be recovered from the measurement
itself. Our contribution is:

\begin{enumerate}
\item A method that segments a raw check-up time series and labels its capacity,
pulse and quasi-OCV sections using only the measured current, voltage and time,
plus a handful of physical scalars that are printed on any cell datasheet or
fixed by the test plan (Section~\ref{sec:method}).
\item A measurement of how much of that method's boundary detection actually
depends on protocol metadata, by ablating the metadata columns and quantifying
what breaks (Section~\ref{sec:res-metadata}). This is a claim that is easy to
assert and rarely tested; we test it.
\item A leave-one-chemistry-out transfer study over NMC, LFP and sodium-ion
cells (Section~\ref{sec:res-loco}), including a negative result --- a
conventionally-trained classifier that fails on an unseen chemistry --- and the
design property that fixes it.
\end{enumerate}

We are explicit about what is \emph{not} claimed. The method is not
metadata-free: it requires physical cell and test parameters, listed in
Section~\ref{sec:inputs}. The distinction we defend is between roughly five
known scalars and a full procedural description of the experiment; we argue in
Section~\ref{sec:discussion} that this distinction is the one that governs
portability.

%% ===================================================================
\section{Related work}
\label{sec:related}

\todobox{Section to write}{%
Cover, in this order: (i) public ageing datasets and their RPT content, read
from the \emph{original} protocol documentation rather than from secondary
processing scripts --- Pozzato and Onori \cite{pozzato}, HNEI, CALCE, Sandia
\cite{snl}, Severson et al.\ \cite{severson}; (ii) BatteryLife
\cite{batterylife} and PyProBE \cite{pyprobe} as the two foils, framed on the
metadata axis of Section~\ref{sec:intro}, not as competitors on accuracy;
(iii) prior work on unsupervised segmentation of battery time series and on
protocol-agnostic feature extraction. Keep the framing on \emph{finding}
structure --- neither foil is claiming to do what this method does, and saying
so plainly is stronger than implying a deficiency.}

%% ===================================================================
\section{Method}
\label{sec:method}

\subsection{Problem statement and inputs}
\label{sec:inputs}

The input is a raw cycler export of one cell: a time series of current $I(t)$,
terminal voltage $V(t)$ and, where available, surface temperature $T(t)$,
covering an arbitrary mixture of ageing and check-up activity. The output is a
partition of that series into segments, each carrying one of the labels CAP,
PUL, qOCV or ``other''.

Besides the signal, the method requires the parameters in
Table~\ref{tab:inputs}. These are physical properties of the cell and targets of
the test plan. They are known before any measurement is made, they are the same
for every cell of a given type, and none of them describes the \emph{order or
composition} of the protocol.

\begin{table}[t]
\centering
\caption{Required inputs besides the measured signal. None describes the
sequence, composition or step indexing of the test.}
\label{tab:inputs}
\begin{tabular}{lll}
\toprule
Parameter & Meaning & Source \\
\midrule
$V_{\max}, V_{\min}$ & voltage limits & datasheet \\
$V_{\mathrm{nom}}$ & nominal voltage & datasheet \\
$Q_{\mathrm{nom}}$ & nominal capacity & datasheet \\
$C_{\mathrm{cap}}$ & capacity-test C-rate & test plan \\
$C_{\mathrm{qocv}}$ & quasi-OCV C-rate & test plan \\
\midrule
$\tau_{\mathrm{pau}}$ & rest duration separating procedures & \todo{state how set} \\
$n_{\min}$ & minimum samples for a segment & \todo{state how set} \\
\bottomrule
\end{tabular}
\end{table}

\todo{$\tau_{\mathrm{pau}}$ and $n_{\min}$ are algorithm parameters rather than
physical ones. State their values, and report the sensitivity of the results to
them --- a reviewer will ask, and an unreported threshold looks tuned.}

\subsection{Segmentation}
\label{sec:segmentation}

Segmentation partitions the series into procedures. A boundary is placed where
the cell's electrical \emph{state} changes --- rest, charge or discharge --- and
where a rest exceeds $\tau_{\mathrm{pau}}$. The state is read from the current,
\begin{equation}
s(t) =
\begin{cases}
\text{rest} & |I(t)| < \varepsilon\,Q_{\mathrm{nom}} \\
\text{charge} & I(t) \ge \varepsilon\,Q_{\mathrm{nom}} \\
\text{discharge} & I(t) \le -\varepsilon\,Q_{\mathrm{nom}},
\end{cases}
\label{eq:state}
\end{equation}
with $\varepsilon$ a dimensionless zero-current tolerance.
Section~\ref{sec:res-metadata} shows the result is insensitive to $\varepsilon$
over two and a half orders of magnitude, so it is not a tuned threshold.

Segments shorter than $n_{\min}$ samples are routed to a discard bucket, as are
the interiors of long rests, which are reduced to their first and last sample so
that relaxation voltage and rest duration remain available downstream without
carrying the intervening samples.

\subsection{Segment features}
\label{sec:features}

Each segment is reduced to a feature vector chosen to be invariant to cell size
and, as far as possible, to chemistry. Currents are divided by
$Q_{\mathrm{nom}}$, making them C-rates; voltages are mapped onto
$[0,1]$ by $(V-V_{\min})/(V_{\max}-V_{\min})$. The vector comprises the current
profile ($I_{\text{mean}}$, $I_{\text{std}}$, $I_{\max}$, $I_{\min}$,
$I_{\text{range}}$, $|I|_{\text{mean}}$), the voltage \emph{edges}
($V_{\max}$, $V_{\min}$, $V_{\text{range}}$), duration
($\Delta t$ and $\log(1+\Delta t)$), and one context feature.

Two choices deserve comment. First, voltage \emph{shape} descriptors (mean and
standard deviation of $V$ within the segment) and temperature are deliberately
excluded: they are the most chemistry- and fixture-bound quantities available,
and omitting them costs no accuracy \todo{cite the ablation number once
Section~\ref{sec:res-ablation} is run}. Second, a segment is not interpretable
in isolation --- a capacity discharge and a preparatory discharge differ mainly
in what preceded them --- so we add the normalised end-of-segment voltage of the
most recent non-rest predecessor. This is the signal a rule-based capacity test
keys on (a true capacity discharge follows a full charge), made available to a
learned model.

\subsection{Segment labelling}
\label{sec:labelling}

We evaluate two labelling routes over identical features.

\paragraph{Unsupervised.} Two-layer HDBSCAN followed by physically-motivated
masks: a capacity segment is a full discharge at $C_{\mathrm{cap}}$ preceded by
a full charge; a quasi-OCV segment is a full sweep at $C_{\mathrm{qocv}}$; a
pulse is a short excursion. The masks are expressed in the normalised feature
space, so they are stated once and evaluated per cell with that cell's
parameters.

\paragraph{Supervised.} A random forest over the same feature vector. The label
space matters more than the model: Section~\ref{sec:res-loco} shows that
training on absolute pipeline labels bakes in the training chemistry's C-rate
and does not transfer, whereas training on a \emph{scale-free} description of
segment shape (full or partial, charge or discharge, at a measured C-rate) and
resolving that description against the target cell's $C_{\mathrm{cap}}$ and
$C_{\mathrm{qocv}}$ at inference time does transfer.

\subsection{Downstream extraction}

Labelled segments feed the quantities a check-up exists to produce: capacity by
coulomb counting the CAP segment, pulse resistance from the PUL segments and
their adjacent rests, and the OCV curve from the qOCV branches. These are
standard and are not a contribution of this work; they are described only so far
as they are used to validate labels in Section~\ref{sec:results}.

%% ===================================================================
\section{Experimental data}
\label{sec:data}

\begin{table}[t]
\centering
\caption{Cells used. All measured at the same institute on the same cycler
family; see Section~\ref{sec:discussion} for what this does and does not
support.}
\label{tab:cells}
\begin{tabular}{llccc}
\toprule
Chemistry & Cell & $Q_{\mathrm{nom}}$ (\si{\ampere\hour}) & $C_{\mathrm{cap}}$ & Cells \\
\midrule
NMC & Sony/Murata VTC6 18650 & 3.0 & C/2 & 7 \\
LFP & A123 APR18650M1B & 1.2 & 1C & 2 \\
Na-ion & \todo{Hina, model} & 10.5 & C/2 & 1 \\
\bottomrule
\end{tabular}
\end{table}

The three chemistries were chosen for what they stress. LFP is the hard case for
any voltage-based feature: its flat plateau compresses the voltage range that
distinguishes a full sweep from a partial one. Sodium-ion is included because
automated check-up decomposition for this chemistry is, to our knowledge,
unreported, and because its differing voltage window tests whether normalisation
by $(V_{\max}-V_{\min})$ is doing real work. The capacity C-rate also differs
between NMC (C/2) and LFP (1C), which turns out to be the decisive variable in
Section~\ref{sec:res-loco}.

\todo{The larger fleet (\num{200} LFP, \num{80} sodium-ion cells) exists but was
not used here. Either bring it in --- it turns per-chemistry $n$ from 1--2 into a
population and makes the transfer result far harder to dismiss --- or state
plainly why a subset was used. A reviewer will notice $n=1$ for sodium-ion.}

%% ===================================================================
\section{Results}
\label{sec:results}

\subsection{How much does segmentation depend on protocol metadata?}
\label{sec:res-metadata}

The claim that structure is recovered from the signal is only as good as the
segmentation stage, which in our implementation reads two columns written by the
cycler: a procedure name and a state label. We therefore ablated them.

Three configurations were compared over the full fleet
(\num{10} cells, \num{131442261} samples, \num{13025} segments):
\textbf{L0}, the production segmenter, using rest-based boundaries, a
procedure-name change, and a state change inside a nominated procedure;
\textbf{L1}, L0 without the procedure-name rule; and \textbf{L2}, in which the
state label is replaced by Eq.~\eqref{eq:state} and the state-change rule is
applied everywhere --- so that \emph{no} protocol column is read. L2 is not
simply ``L1 with more removed'': it also promotes the state rule to a general
one. That is deliberate. The question is not how much a stripped-down segmenter
loses, but whether the signal carries the structure the protocol columns supply.

Because the ablation re-implements the segmenter, its L0 output is verified
against the production implementation sample by sample; a single differing
sample aborts the comparison. All ten cells passed.

\paragraph{The state column carries no information.}
Comparing the cycler's state label against Eq.~\eqref{eq:state}, \num{1100} of
\num{131442261} samples disagree (\SI{0.0008}{\percent}). Recall of the rest
class is exactly \num{1.000} on every cell; precision is at least
\num{0.99986}, and is flat for $\varepsilon$ between \num{e-4} and \num{3e-2}
--- a factor of \num{300}. The residual disagreement is not a chemistry effect
but a few hundred samples of cycler bookkeeping states recorded at zero current,
plus roughly \num{600} samples caught at a current zero crossing.

\paragraph{The procedure name fires often and decides little.}
Table~\ref{tab:rule3} shows that a majority of procedure-name boundaries occur
where no other rule fires --- which is why the dependence looked severe --- but
that \SIrange{81}{87}{\percent} of those cut inside a long rest or a
sub-threshold fragment that is discarded regardless. The naive count overstates
the dependence by roughly a factor of six.

\begin{table}[t]
\centering
\caption{Procedure-name boundaries. ``Unique'' counts firings where no other
rule fires; ``discarded'' counts the subset of those that cut only material the
segmenter throws away.}
\label{tab:rule3}
\begin{tabular}{lrrr}
\toprule
Chemistry & Firings & Unique & Discarded \\
\midrule
NMC    & \num{2177} & \num{1242} (\SI{57.1}{\percent}) & \num{1037} (\SI{83.5}{\percent}) \\
LFP    & \num{585}  & \num{358} (\SI{61.2}{\percent})  & \num{312} (\SI{87.2}{\percent}) \\
Na-ion & \num{78}   & \num{16} (\SI{20.5}{\percent})   & \num{13} (\SI{81.2}{\percent}) \\
\bottomrule
\end{tabular}
\end{table}

\paragraph{Check-up segments survive without any protocol column.}
Table~\ref{tab:checkup} gives the fate of every segment that the pipeline
ultimately labels CAP, PUL or qOCV. We distinguish a segment that is
\emph{split} --- still present, cut into pieces --- from one that is
\emph{merged} into a neighbour or \emph{lost}, which cannot be recovered
downstream. At L2 there are no merges and no losses on any chemistry: every
check-up segment survives intact.

\begin{table}[t]
\centering
\caption{Fate of check-up segments (CAP, PUL, qOCV) relative to the L0
reference. \emph{Intact} = exact or split.}
\label{tab:checkup}
\begin{tabular}{lrrrrr}
\toprule
& & \multicolumn{2}{c}{L1 (no procedure name)} & \multicolumn{2}{c}{L2 (no protocol column)} \\
\cmidrule(lr){3-4}\cmidrule(lr){5-6}
Chemistry & $n$ & exact & intact & exact & intact \\
\midrule
NMC    & \num{3054} & \SI{89.3}{\percent} & \SI{89.3}{\percent} & \SI{99.9}{\percent} & \SI{100}{\percent} \\
LFP    & \num{332}  & \SI{85.5}{\percent} & \SI{85.5}{\percent} & \SI{99.7}{\percent} & \SI{100}{\percent} \\
Na-ion & \num{90}   & \SI{100}{\percent}  & \SI{100}{\percent}  & \SI{95.6}{\percent} & \SI{100}{\percent} \\
\bottomrule
\end{tabular}
\end{table}

The shape of the L1 loss is diagnostic. Every one of the \num{376} segments L1
merges is a pulse (\num{349} test pulses, \num{27} restore pulses); capacity and
quasi-OCV segments never break at any level. Inspecting the procedure-name
boundaries that L1 loses and L2 recovers shows why: they are rest-to-charge
transitions following a rest \emph{shorter} than $\tau_{\mathrm{pau}}$. No
rest-based rule can fire there, so L1 has nothing to cut with --- but the state
changes, and that change is present in the current. The procedure name was
supplying timing, not information.

That L2 outperforms L1 while reading strictly less metadata is the central
result of this section: the boundaries are in the signal, and the protocol
columns were a convenience.

\paragraph{What L2 does cost.} Outside the check-up taxonomy, L2 disturbs the
rest structure substantially (on NMC, \num{1161} merged and \num{302} lost
non-check-up segments). Rest stubs supply the predecessor-voltage context
feature of Section~\ref{sec:features} and the relaxation window used for pulse
resistance, so a segmenter built on L2 would need to treat rests explicitly.
Separately, impedance-spectroscopy dwell windows are indistinguishable from
rests by current alone. Neither affects the CAP/PUL/qOCV result, and both are
stated here rather than deferred to the limitations section because they bound
what the result licenses.

\subsection{Cross-chemistry transfer}
\label{sec:res-loco}

We next ask whether the \emph{labelling} stage transfers to a chemistry it has
never seen, given only that chemistry's physical parameters and no labelled
examples from it.

The unsupervised route transfers directly. Applied to LFP with only
Table~\ref{tab:inputs} changed, it recovers all \num{16} check-ups, yielding a
monotonic capacity trajectory from \SI{98.3}{\percent} to \SI{60.2}{\percent} of
nominal, with a first measured capacity of \SI{1.18}{\ampere\hour} against
\SI{1.2}{\ampere\hour} nominal. The mechanism is that the capacity C-rate enters
as a parameter rather than being learned: on LFP the first clustering layer
finds no clean capacity cluster --- the flat plateau removes the voltage
contrast it relies on --- and a second, stricter layer recovers it.

The supervised route exposes a design failure worth reporting. A random forest
trained on NMC and applied to LFP labels \emph{none} of the \num{16} capacity
discharges correctly, and every check-up consequently reports no state of
health. The cause is not the flat plateau but the C-rate: NMC capacity tests run
at C/2 and LFP at 1C, and a model trained on absolute pipeline labels has
learned ``capacity test'' to mean ``a full discharge at C/2''. Retraining the
same model on a scale-free description of segment shape --- full or partial,
charge or discharge, at a measured C-rate --- and resolving that description
against the target cell's $C_{\mathrm{cap}}$ at inference recovers all \num{16},
reproducing the unsupervised result exactly. On sodium-ion, both routes recover
all \num{6} check-ups.

\todobox{Provenance and completeness}{%
These numbers are from a run of 2026-07-23 whose comparison artefacts have since
been overwritten (the comparison tool writes to a fixed filename, so a later
chemistry's run replaced the earlier one's output). \textbf{Re-run all three
chemistries and archive the outputs under per-chemistry names before this
section is submitted.} Also report, in the table this paragraph should become:
per-chemistry counts of segments and check-ups, precision and recall per label
class rather than capacity counts alone, and the sodium-ion over-tagging
behaviour --- the scale-free model tags two capacity segments per check-up in
four of six programmes, because it resolves capacity by C-rate alone and does
not test for a full voltage swing. The downstream capacity calculation rejects
the extras, so the reported capacity is unaffected, but the label precision is
not \SI{100}{\percent} and the paper must say so.}

\subsection{Segment-type accuracy against hand labels}
\label{sec:res-f1}

\todobox{Experiment not yet run}{%
Hand-label a stratified sample of segments per chemistry (target: a few hundred
each, drawn across cells and ageing states, labelled by someone who did not
build the method) and report per-class precision, recall and F1 for CAP, PUL,
qOCV and other, for both labelling routes. This is the only place in the paper
where the labels are compared against ground truth rather than against each
other --- Section~\ref{sec:res-loco} compares two automated routes, which cannot
detect an error both make. Without it, a reviewer is entitled to ask what the
accuracy actually is.}

\subsection{Check-up detection recall against the test plan}
\label{sec:res-recall}

\todobox{Experiment not yet run}{%
The test plan states when every check-up was scheduled, which is ground truth
that requires no labelling effort: for each cell, compare the set of detected
check-ups against the scheduled set and report recall and false-detection rate.
This measures the question the paper is actually about --- \emph{did we find the
check-ups} --- as opposed to whether we labelled their contents correctly.}

\subsection{Comparison against a rule-based baseline}
\label{sec:res-baseline}

\todobox{Experiment not yet run}{%
Report the method against the fixed-threshold approach the foils use, applied to
our data: a hard current threshold for capacity detection of the kind used in
\cite{batterylife}. Report where it fails and why. The comparison must be fair
--- the baseline was written for a different purpose --- so frame it as
quantifying the cost of a fixed threshold under a chemistry change, not as a
defeat of prior work.}

\subsection{Feature ablation}
\label{sec:res-ablation}

\todobox{Experiment not yet run}{%
Two ablations are promised earlier in the text and must be run:
(i) adding back voltage-shape and temperature features, to support the claim in
Section~\ref{sec:features} that excluding them costs no accuracy;
(ii) the LFP plateau case --- show the failure mechanism of the
voltage-range feature on a flat OCV curve, then the coulombic replacement
($|I|\,\Delta t / Q_{\mathrm{nom}}$, the charge swept as a fraction of nominal
capacity, already implemented) and the recovery it produces. Mechanism, failure,
fix, number.}

\subsection{Validation on an external dataset}
\label{sec:res-external}

\todobox{Experiment not yet run}{%
All cells in Table~\ref{tab:cells} come from one institute, so nothing here
supports a cross-laboratory claim. The strongest available remedy is the
Pozzato--Onori dataset \cite{pozzato}, which contains the full taxonomy ---
C/20 capacity, HPPC pulses and pseudo-OCV --- from an independent laboratory.
Ingest it with a thin format adapter (which itself demonstrates the thesis:
adapter, not analysis code) and report the same metrics.
\textbf{Verify the file structure against the original dataset, not against a
secondary processing script.} If this cannot be completed, the portability
claim must be explicitly scoped to within-laboratory in the abstract and the
conclusion, not only in the limitations.}

%% ===================================================================
\section{Discussion}
\label{sec:discussion}

\subsection{What is and is not claimed}

The method is not metadata-free, and describing it that way would be wrong. It
requires the parameters of Table~\ref{tab:inputs}. What it does not require is a
description of the experiment's structure: which step is which, in what order,
under what indexing. That is the distinction that governs portability, because
the two kinds of information have different availability. Voltage limits,
nominal capacity and target C-rates are known before the cell is connected, are
identical for every cell of a type, and are stable across laboratories. A
step-index map is authored per protocol, is invalidated when the protocol is
edited, and does not exist at all for a dataset one did not generate.

\subsection{Where portability comes from}

Section~\ref{sec:res-loco} isolates the mechanism. The two routes that transfer
--- the rule-based one and the scale-free supervised one --- share the property
that the chemistry-specific quantity ($C_{\mathrm{cap}}$) enters as a
\emph{parameter at inference time} rather than being absorbed into learned
weights. The route that fails absorbed it. This is a statement about how the
label space is designed, not about model capacity, and it suggests that
portability across chemistries is achieved by choosing what to parameterise
rather than by training on more chemistries.

\subsection{Limitations}

\begin{itemize}
\item \textbf{Single laboratory.} All cells in Table~\ref{tab:cells} were
measured at one institute on one cycler family, so cross-laboratory portability
is argued from construction, not demonstrated
(Section~\ref{sec:res-external}).
\item \textbf{The signal-only segmenter was measured, not adopted.}
Section~\ref{sec:res-metadata} evaluates it as an ablation of the shipped
implementation. We report what it recovers, not a system built on it, and the
rest-structure cost noted there would have to be addressed first.
\item \textbf{Check-ups, not cycling.} The method is developed and evaluated on
check-up data. Nothing here establishes behaviour on cycling profiles.
\item \textbf{Small $n$ per chemistry} for LFP and sodium-ion as reported here
\todo{resolve with the larger fleet, see Section~\ref{sec:data}}.
\item \textbf{Two residual metadata uses} remain outside boundary detection: a
state-based filter drops a small class of cycler bookkeeping rows, and files are
grouped by test identifier. Neither describes protocol composition, and the
second has a single-group fallback, but both should be disclosed rather than
defended.
\end{itemize}

%% ===================================================================
\section{Conclusion}
\label{sec:conclusion}

\todobox{Write last}{%
Write this once Sections~\ref{sec:res-f1}--\ref{sec:res-external} carry results.
It should state the two established findings --- that check-up boundaries are
recoverable from the current alone with no loss of capacity, pulse or quasi-OCV
segments across three chemistries, and that cross-chemistry label transfer is a
property of the label space rather than of training breadth --- and should not
reach beyond them.}

%% ===================================================================
\section*{Data and code availability}
\todo{state repository, licence, and which figures/tables regenerate from which
command}

\section*{Acknowledgements}
\todo{funding, project number}

\bibliographystyle{elsarticle-num}
\bibliography{refs}

\end{document}
